%\documentclass[notes,usenames,dvipsnames]{beamer}
%\documentclass[notes]{beamer}       % print frame + notes
%\documentclass[notes=only]{beamer}   % only notes
\documentclass[usenames,dvipsnames]{beamer}             % only frames
\usepackage{pgfpages}
%\setbeameroption{show notes on second screen}
%\setbeameroption{show notes}
%subfigures
\usepackage{caption}
\usepackage{subcaption}

%tables packages
\usepackage{multirow}

% math
\usepackage{amsmath}

% bash command
\usepackage{graphicx}
\usepackage{listings}
% varbatim for ascii figures
\usepackage[T1]{fontenc}
\usepackage[utf8]{inputenc}
%\usepackage{verbatim}
\usepackage{lipsum} % for context
\usepackage{fancyvrb}
\usepackage{varwidth}

% notes prefixed in pympress
\addtobeamertemplate{note page}{}{\thispdfpagelabel{notes:\insertframenumber}}
%theme used
\usetheme{Madrid}
%Information to be included in the title page:
\title[Computer Architecture] %optional
{Arhitectura Calculatoarelor}

\subtitle{Seria AB - Curs 0 - Introducere}

\author[Ștefan-Dan Ciocîrlan] % (optional, for multiple authors)
{}

\institute[NUSTPB] % (optional)
{
  \inst{}%
  National University of Science and Technology\\
  POLITEHNICA Bucharest
}

\date[NUSTPB 2025] % (optional)
{Computer Architecture}

\logo{
\includegraphics[height=1.3cm]{../../media/logo-poli-color9.png}
%\includegraphics[height=0.7cm]{media/sbipLogo.pdf}
}


% Roman numerals
\newcommand*{\rom}[1]{\expandafter\@slowromancap\romannumeral #1@}
%split
\usepackage{amsmath}
%colors
%\usepackage[usenames,dvipsnames]{color} %loaded by the dcoument class
%subfigure
\usepackage{subcaption}
% block over block uncover
\setbeamercovered{invisible}
%\setbeamercovered{transparent}

%extra slide content
\AtBeginSection[]
{
  \begin{frame}
    \frametitle{Conținut}
    \tableofcontents[currentsection]
  \end{frame}
}
%notes or not
%\setbeamertemplate{note page}[plain]

\begin{document}

\frame{\titlepage}


\section{Prezentarea materiei}

\begin{frame}
    \frametitle{Parcurs academic}
    \begin{itemize}
        \item Licență: UPB-ACS/IPP-ENSTA "Interacțiunea om-robot într-un context de asistență socială"
        \item Disertație: UPB-ACS/IPP-ENSTA "Social engineering attack using humanoid robots"
        \item Doctorat: UPB-ACS/NUS-SOC "Number representation systems in computer engineering"
        \item Post-Doctorat: NUS-SOC "Singapore Blockchain Innovation Programme"
        \item Proiecte: CuEVM \footnotemark[1], NRSL/HNRSL, ENRICHME
        \item Interese academice: arhitectura calculatoarelor, matematică, blockchain, criptografie
    \end{itemize}
    \note{
    }
    \footnotetext[1]{https://github.com/sbip-sg/CuEVM}
\end{frame}

\begin{frame}
    \frametitle{Parcurs didactic}
    \begin{itemize}
        \item Instructor Hackademy: CCNA(2014-2016), C++ (2015)
        \item Tutor UPB-ACS (2015-2018): Calculatoare numerice 1 (2015), Calculatoare numerice 2 (2016), Procesarea semnalelor (2017), Rețele locale (2017), Proiectarea rețelelor (2017).
        \item Asistent UPB-ACS (2019-2021): Arhitectura calculatoarelor (2019), Calculatoare numerice 1 (2020), Calculatoare numerice 2 (2020), Arhitectura Sistemelor de calcul (2022),  Rețele locale (2019-2020), Proiectarea rețelelor (2019-2021), Proiectarea rețelelor neurale private(2022), Tehnici de protecție a vieții private (2019-2021).
        \item Asistent NUS-SOC (2022-2023): Cloud Computing (2022)
        \item Guest Lecturer NUS-SOC (2023-2024): Cloud Computing (2023)
        \item Școli de vară : UPB-ACS Arhitectura calculatoarelor (2022), UPB-ACS Data Science (2022), NUS-SOC Cloud Computing (2023)
        \item Interese didactice: open source, corectare automată
    \end{itemize}
    \note{
    }
\end{frame}

\begin{frame}
    \frametitle{Cuprins materie}
    \begin{enumerate}
        \setcounter{enumi}{0}
        \item Introducere
        \item Structură Calculator Numeric
        \item Reprezentarea și prelucrarea informației
        \item Memorii
        \item UAL
        \item Arhitectura Calculatorului Didactic
        \item Limbaje de asamblare
        \item Subsistem intrare/ieșire
        \item Întreruperi
        \item Microprogramare
    \end{enumerate}
\end{frame}
    
\begin{frame}
    \frametitle{Obiectivele materiei}
    
    \begin{itemize}
        \item Cultură generală despre arhitectura calculatoarelor
        \begin{itemize}
            \item Reprezentarea informației
            \item Memorii
            \item Seturi de instrucțiuni
            \item Codificare instrucțiuni
            \item Limbaj de asamblare
            \item Întreruperi de procesor
        \end{itemize}
        \item Descriere hardware a unui procesor prin limbajul Verilog
    \end{itemize}
\end{frame}




\begin{frame}
    \frametitle{Platforma cursului}
    
    \begin{itemize}
        \item \href{https://github.com/cs-pub-ro/computer-architecture}{GitHub: https://github.com/cs-pub-ro/computer-architecture}
        \item \href{https://cs-pub-ro.github.io/computer-architecture/}{Open Education Hub: https://cs-pub-ro.github.io/computer-architecture/}
        \item \href{https://curs.upb.ro/2024/course/view.php?id=1815}{Moodle: https://curs.upb.ro/2024/course/view.php?id=1815}
    \end{itemize}
\end{frame}

\section{Notare}

\begin{frame}
\frametitle{Notare}

\begin{itemize}
    \item 50p examen final
    \item 50p laborator
    \begin{itemize}
        \item 10p activitate laborator (prezență și teste surpriză)
        \item 10p Test circuite combinaționale
        \item 10p Test circuite secvențiale
        \item 20p Colocviu final calculator didactic
    \end{itemize}
    \item 10p Tema (bonus)
\end{itemize}
\end{frame}

\begin{frame}
\frametitle{Promovare}
    
\begin{itemize}
    \item UNSTPB/NUSTPB
    \begin{itemize}
        \item Minim 25p Laborator
        \item Minim 50p Total
    \end{itemize}
\end{itemize}
\end{frame}

\begin{frame}
\frametitle{Test circuite combinaționale}

\begin{itemize}
    \item Materia necesară: laboratoarele 1, 2 și 3
    \item Durată: 60 de minute
    \item Când: la începutul laboratorului 4
    \item Structură: 3 exerciții practice de implementat în Verilog
    \item Locație: sala de laborator
    \item Platformă: Calculator Laborator - Safe Exam Browser - Moodle - VPL
\end{itemize}
\end{frame}

\begin{frame}
\frametitle{Test circuite secvențiale}

\begin{itemize}
    \item Materia necesară: laboratoarele 4, 5 și 6
    \item Durată: 60 de minute
    \item Când: la începutul laboratorului 7
    \item Structură: 3 exerciții practice de implementat în Verilog
    \item Locație: sala de laborator
    \item Platformă: Calculator Laborator - Safe Exam Browser - Moodle - VPL
\end{itemize}
\end{frame}

\begin{frame}
\frametitle{Colocviu final calculator didactic}

\begin{itemize}
    \item Materia necesară: toate laboratoarele
    \item Durată: 120 de minute
    \item Când: la ultimul laborator
    \item Structură:
    \begin{itemize}
        \item 3 exerciții practice de implementat în Verilog
        \item 20 întrebări teoretice
    \end{itemize}
    \item Locație: sala de laborator
    \item Platformă: Calculator Laborator - Safe Exam Browser - Moodle - VPL/Quiz
\end{itemize}
\end{frame}

\begin{frame}
\frametitle{Tema}

\begin{itemize}
    \item Durată: Întreg semestrul cu termen limită începutul ultimii săptămâni de laborator
    \item Platformă: Moodle - VPL
    \item Punctaj: Se va adăuga punctajului de laborator (nu se ia în considerare la punctajul minim)
    \item Punctaj: Punctajul total de laborator se trunchează la 50 de puncte.
\end{itemize}
\end{frame}

\begin{frame}
\frametitle{Examen Final (AB)}

\begin{itemize}
    \item Platformă: Moodle - VPL/Quiz
    \item Durată: 180 de minute
    \item Structură:
    \begin{itemize}
        \item 20 de întrebări teoretice
        \item 1 problemă
    \end{itemize}
\end{itemize}
\end{frame}


\section{Reguli}
\begin{frame}
    \frametitle{Rolurile cadrelor didactice (părere personală)}
\begin{itemize}
    \item De a oferi know-how-ul necesar pentru a vă putea dezvolta în domeniu.
    \item De a fi onest și predictibil.
    \item De a oferi resursele necesare.
    \item De a fi responsabil pentru un mediul sigur.
\end{itemize}

\end{frame}

\begin{frame}
    \frametitle{Implicații}
\begin{itemize}
    \item Nu veți avea întrebări de reproducere a informației.
    \item Nu e un obiectiv să aveți note mici, ci să existe o diferențiere între voi.
    \item Aveți model pentru a vă pregăti pentru examen.
    \item Nu voi accepta comportament abuziv.
\end{itemize}

\end{frame}

\begin{frame}
    \frametitle{Laborator}
    
\begin{itemize}
    \item Schimbarea intervalului asignat în orar este posibilă până la al doilea laborator, dar numai prin schimb cu altă persoană și cu acordul asistenților responsabili de intervalele implicate (nu se pot face schimbări între serii).
    \item Recuperarea unui laborator (pierdut sau în avans) este posibilă doar în limita locurilor disponibile la calculatoarele din sală. Înainte de a vă prezenta pentru recuperarea unui laborator trimiteți un e-mail asistentului vostru și asistentului din intervalul în care vreți să recuperați, pentru a obține acordul lor.
\end{itemize}
\end{frame}


\begin{frame}
    \frametitle{Regulament}
    
    \href{https://cs-pub-ro.github.io/computer-architecture/rules}{https://cs-pub-ro.github.io/computer-architecture/rules}

    Putem aveam "gentlemen's agreement". Propuneri pot veni din partea voastră.

    \begin{itemize}
        \item Prezență cursuri din partea companiilor AMIQ 08.10.2025, NXP 12.11.2025, TODO: Onsemi, ChipAhead, Infineon
    \end{itemize}

\end{frame}

\begin{frame}
    \frametitle{Feedback 2024-2025}
    
    \begin{itemize}
        \item Prea multe platforme - Moodle, GitHub, Teams -> Github/Moodle
        \item Asistenții nu știu Verilog -> RHDL
        \item Asistentul Grinch -> Nu mai este asistent
        \item Testele VPL prea grele -> Proiect
        \item Geoguessr -> Păstrăm
        \item Proful nu ia în serios materia -> Aducem invitați din industrie
    \end{itemize}
\end{frame}

\section{Q\&A}
\begin{frame}
\end{frame}

%\begin{frame}
%\frametitle{Table of Contents}
%\tableofcontents
%\end{frame}


\end{document}